\section{古典データの符号化}
\label{sec:data-encoding}

古典データを量子アルゴリズムで扱うには，データを量子状態に変える必要がある．
この変換を符号化という．
本論文では，基底符号化と振幅符号化の二つを用いる\cite{shimada2020}．

古典的な$n$ビット列$x=x_{n-1}\cdots x_0$を$n$量子ビットの計算基底状態
$\ket{x_{n-1}\cdots x_0}$へ対応させる方法を基底符号化という．
この方法は，位相推定で整数や小数をレジスタに保存するときに用いる．

\begin{defi}[振幅符号化]
非零ベクトル$\bm b=(b_0,\ldots,b_{N-1})^{\mathsf T}\in\mathbb C^N$に対し，
\begin{equation}
 \ket b=\frac1{\lVert\bm b\rVert_2}
          \sum_{j=0}^{N-1}b_j\ket j
 \label{eq:amplitude-encoding}
\end{equation}
を$\bm b$の振幅符号化という．次元$N$が2の冪でない場合は，
零成分を加えて$2^n$次元の状態空間へ埋め込む．
\end{defi}

\begin{prop}
式\eqref{eq:amplitude-encoding}で定まる$\ket b$は規格化されている．
\end{prop}
\begin{proof}
計算基底の正規直交性を用いると
\[
 \braket{b|b}
 =\frac1{\lVert\bm b\rVert_2^2}
   \sum_{j,k}\overline{b_j}b_k\braket{j|k}
 =\frac{\sum_j|b_j|^2}{\lVert\bm b\rVert_2^2}=1
\]
である．非零性の仮定$\bm b\ne\bm0$により分母は零でない．
\end{proof}

HHLへの入力は式\eqref{eq:amplitude-encoding}の$\ket b$である．
ただし，古典的なベクトル$\bm b$が与えられていても，
量子状態$\ket b$を短い時間で準備できるとは限らない．
そのため，第4章では状態を準備する方法と必要な計算量を条件として書く．
